\chapter{Om krydsreferencer}

Hvis ikke man kender til de rige muligheder for automatisk håndtering
af referencer, som \latex har, er det som at trække rundt på
cykelstien med en dyr racercykel og ærgre sig over, at man ikke kan
komme hurtigere frem, fordi man er nødt til at trække rundt med hvad
man tror er et dårligt designet løbehjul. 

\section{Krydsreferencer} \label{afs:kryds}

Afsnit, underafsnit, tabeller, figurer og formler kan få numre i
\latex. Det gør man med \verb2\label2-kommandoe, som placerer et
mærke. Det kan man senere referere til med en \verb2\ref2. Se efter i
kildeteksten til dette dokument for at lære, hvordan man laver en
henvisning til afsnit \ref{afs:kilder}. 

Det er aldrig nok at skrive \verb2\ref2, for denne kommando skriver
blot nummeret på den label, man henviser til. Jeg har defineret en
kommando \verb2\afs2, der tager et mærke som argument og skriver
"Afsnit'' og nummeret på mærket. Se, hvordan denne kommando er
defineret og brugt. Jeg bruger den også her i \afs{afs:kryds}. 

Pakken \verb2cleveref2 tilbyder en lignende funktionalitet og er værd
at overveje at bruge. 

\section{Kildehenvisninger} \label{afs:kilder}

\begin{ovelse}
  Prøv at skrive
\begin{verbatim}
\bibliographystyle{apalike}
\end{verbatim}
  i stedet for \verb2 \bibliographystyle{plain}2 og se, hvad der
  sker. Hvis du ikke bruger Overleaf, så husk at kompilere to gange. 
\end{ovelse}

%%% Local Variables:
%%% mode: latex
%%% TeX-master: "report"
%%% End:
